\subsection{Excluding \texorpdfstring{$C_4\times C_{16}$}{C4 x C16}}
\begin{proposition}\label{prop:decic416}
No elliptic curve over a decic number field contains $C_4\times C_{16}$.
\end{proposition}
\begin{proof}
Put $F=\Q(i)$, $X=X_1(4,16)$ and $J=\operatorname{Jac}(X)$.
A proposed subgroup gives a noncuspidal effective $F$-rational
divisor $D$ of degree $5$. The curve has genus $13$, and
$J(F)$ has rank zero by \cite[Theorem~4.1]{DS}.
Its $F$-gonality is at least $8$ by \cite[Table~1]{NV}.
Lemma~\ref{lem:decic-cuspidal-reduction} excludes entirely cuspidal
reduction at the primes above $3$ and $5$.

Let $C_F$ be generated by the differences of the global cuspidal
orbit divisors, normalized at a rational cusp $c_0$.
The $24$ orbit divisors have degrees $1^{16},2^4,4^4$.
Their integral modular-unit valuation vectors identify the same
divisors in both reductions. Over $\mathbb F_5$ we compute
\[
\begin{aligned}
 J(\mathbb F_5)&\simeq[2,2,2,2,2,2,4,4,40,40,80,80],\\
 C_F&\simeq[2,2,2,4,4,20,40,40,80].
\end{aligned}
\]
The quotient of the first group by the cusp image is $C_2^5$.
Odd-prime torsion injectivity therefore gives $2J(F)\subseteq C_F$.
We do not assume that the cusps generate all of $J(F)$.

The noncuspidal counts in degrees one through five are
\[
\begin{array}{c|rrrrr}
 &1&2&3&4&5\\\hline
 \mathbb F_5&0&0&64&160&512\\
 \mathbb F_9&0&0&320&1808&11904.
\end{array}
\]
They agree with the Frobenius polynomials of $J_H(64)$,
$H=\langle17\rangle$, using squared Frobenius at $3$.
At $5$ the cusp signature is $1^{16},2^4,4^4$, so the number of
noncuspidal degree-$5$ patterns is
\[
 64\left(\binom{17}{2}+4\right)+160\cdot16+512=12032.
\]
We include every one of these divisors.

At $9$ the cusp signature is $1^{24},2^8$.
The $64$ pairing-preserving level automorphisms partition the
cubic, quartic and quintic supports into $5$, $33$ and $186$ orbits.
We include every cusp complement after normalizing the noncuspidal
support. There are
\[
 5\left(\binom{25}{2}+8\right)+33\cdot24+186=2518
\]
target representatives.
This normalization is valid because the full source set at $5$
is stable under the global automorphisms. The doubled degree-$5$
class transforms affinely, with cusp translation
$10[g(c_0)-c_0]$, and the actual common cusp maps account for it.
The finite-field embeddings fix the chosen image of $i$.

The common cusp maps have equal kernels. Comparing the two images
of $2[D-5c_0]$ leaves $158$ target classes: $157$ cubic-plus-cusp
patterns and one quintic pattern. A separate integer-lattice replay
reconstructs the exact common graph, checks its order $327680000$,
and reproduces all matches from the actual relation matrices.

For each remaining class we compute a signed cusp vector $v$ with
\[
 2D\sim L_v:=10c_0+\sum_{j=1}^{24}v_j(C_j-\deg(C_j)c_0).
\]
The equality lifts globally by odd-prime torsion injectivity.
At a prime above $13$, the cusp signatures agree with the same
global divisors. We compute the complete systems
$|\overline L_v|(\mathbb F_{13})$.
For $150$ classes the Riemann--Roch dimension is one, and for eight
it is two. We therefore test
\[
 150+8(13+1)=262
\]
rational members. Every member has an odd support multiplicity.
However, a specialization of $2D$ has only even multiplicities.
This excludes all $158$ classes and proves the proposition.

The independent replay checks every signed cusp preimage, all
$12032$ source divisors, the $2518$ target representatives and their
full orbit coverage, the $158$ diagonal matches, and all $262$
parity decisions. The finite class groups, divisor class maps and
Riemann--Roch calculations use proved native Magma arithmetic.
\end{proof}
